\section{Welfare and continuation selection}\label{sec:welfare}

The equilibrium multiplicity in Proposition \ref{prop:equilibrium} matters for the welfare comparison because markets are segmented and firms can set country-specific prices. The price terms cancel under a common symmetric continuation across the two union markets, but they need not cancel when the two markets select different continuation prices. This section makes that scope explicit.

Let the symmetric member price in union market A be $s_A$ and that in union market B be $s_B$. In either foreclosed market, the two active member firms each sell $3/2$ units. With equal prices within a market, each active firm serves a half-arc of length $1/2$ and a half-arc of length $1$, so total transportation cost in either union market is
\begin{align}\label{eq:transport-foreclosure}
 T^{\SU}
 &=2\int_{0}^{1/2}x^2\,dx+2\int_{0}^{1}x^2\,dx
 =\frac34.
\end{align}
Consumer surplus in country A is therefore
\begin{equation}\label{eq:cs-su}
 CS_A^{\SU}=3V-3s_A-\frac34.
\end{equation}

National welfare includes the domestic firm's worldwide profit. Firm A sells $3/2$ units in its home union market at price $s_A$ and $3/2$ units in the other union market at price $s_B$. Under the recognition-stage behavior used in the original proof of Proposition 3, the nonmember country recognizes all standards, and the member firm earns one additional unit of profit there \citep[p.~382]{gandalshy2001}. Hence
\begin{equation}\label{eq:profit-su}
 \Pi_A^{\SU}=\frac32 s_A+\frac32 s_B+1.
\end{equation}

\begin{proposition}[Welfare under continuation selection]\label{prop:welfare}
Suppose the two union markets use symmetric foreclosed pure-strategy continuations with member prices $s_A$ and $s_B$, and retain the nonmember recognition behavior used in Proposition 3 of \citet{gandalshy2001}. Then country A's welfare is
\begin{equation}\label{eq:ts-crossmarket}
 TS_A^{\SU}
 =3V+\frac14+\frac32(s_B-s_A),
\end{equation}
with the symmetric expression for country B.

If the two union markets select a common symmetric member price $s_A=s_B=s$, then
\begin{equation}\label{eq:ts-su}
 TS_M^{\SU}=3V+\frac14.
\end{equation}
Under mutual recognition,
\begin{equation}\label{eq:ts-mr}
 TS^{\MR}=3V-\frac14,
\end{equation}
so the member-country gap under a common symmetric continuation is
\begin{equation}\label{eq:welfare-gap}
 TS_M^{\SU}-TS^{\MR}=\frac12>0.
\end{equation}
Thus the original member-country ranking is preserved under a common symmetric continuation, including the symmetric continuation selected within the cost-floor equilibrium class of Proposition \ref{prop:equilibrium}. In the unrestricted original game, however, the welfare comparison is not invariant to arbitrary market-by-market continuation selection.
\end{proposition}

\begin{proof}
Adding \eqref{eq:cs-su} and \eqref{eq:profit-su} gives
\[
 TS_A^{\SU}
 =3V-3s_A-\frac34+\frac32s_A+\frac32s_B+1
 =3V+\frac14+\frac32(s_B-s_A),
\]
which proves \eqref{eq:ts-crossmarket}. Setting $s_A=s_B$ eliminates the cross-market transfer term and gives \eqref{eq:ts-su}.

Under mutual recognition all standards are accepted in every country and all three firms are symmetric. In a representative market each firm charges $1$ and sells one unit; consumer expenditure is $3$ and transportation cost is $1/4$, so consumer surplus is $3V-13/4$. The domestic firm earns one unit in each of the three segmented markets, giving worldwide profit $3$. Hence $TS^{\MR}=3V-1/4$, and subtraction yields \eqref{eq:welfare-gap}.
\end{proof}

The welfare result therefore has two layers. Under a common symmetric continuation, prices redistribute surplus between consumers and the domestic producer but cancel from national welfare, reproducing the original $1/2$ member-country gap. Under unrestricted market-specific continuation selection, that cancellation is incomplete because the price paid by country A's consumers is $s_A$ while its domestic firm also earns export profit at $s_B$. Equation \eqref{eq:ts-crossmarket} records this continuation-selection dependence. It does not by itself establish a government-stage policy reversal, because the government-stage equilibrium and equilibrium-selection rule under the unrestricted continuation correspondence are not characterized here.
